% Source: https://artofproblemsolving.com/community/c3h1565878_complex_numbers
% Stable identity: aops-topic-1565878-q01
Multiply both sides by $z+1$ (note $z\ne-1$): $z=(-1+i)(z+1)=(-1+i)+(-1+i)z$.

Now collect the $z$ terms on the left by subtracting $(-1+i)z$ from both sides:
\[
z-(-1+i)z=-1+i\qquad\Longrightarrow\qquad\bigl(1-(-1+i)\bigr)z=-1+i\qquad\Longrightarrow\qquad(2-i)z=-1+i,
\]
because $1-(-1+i)=1+1-i=2-i$. Dividing, and multiplying numerator and denominator by the conjugate $2+i$:
\[
z=\dfrac{-1+i}{2-i}=\dfrac{(-1+i)(2+i)}{(2-i)(2+i)}=\dfrac{-2-i+2i+i^2}{5}=\dfrac{-3+i}{5}.
\]
So $z=\boxed{-\dfrac{3}{5}+\dfrac{1}{5}i}$.

Check: $z+1=\dfrac{2+i}{5}$ and $\dfrac{z}{z+1}=\dfrac{-3+i}{2+i}=\dfrac{(-3+i)(2-i)}{5}=\dfrac{-6+3i+2i+1}{5}=-1+i$.
